\subsection{Intersecciones completas relativas (caso plano).}
\label{subsection:IV.19.3}

\oldpage[IV-4]{194}

\begin{definition}[19.3.1]
\label{IV.19.3.1}
Sea $A$ un anillo local noetheriano. Se dice que $A$ es un \emph{anillo de intersección completa} (o también un \emph{anillo de intersección completa absoluta}, si hay riesgo de confusión) si la completación $\widehat A$ es isomorfa al cociente de un anillo local noetheriano completo regular $B$ por un ideal regular (es decir, (16.9.7), un ideal engendrado por una sucesión regular de elementos de $B$).

Todo anillo local regular es un anillo de intersección completa.

Se dice que un preesquema localmente noetheriano $X$ es una \emph{intersección completa} (absoluta) \emph{en el punto} $x\in X$ si el anillo $\sh{O}_{X,x}$ es un anillo de intersección completa.
\end{definition}

\begin{proposition}[19.3.2]
\label{IV.19.3.2}
Sean $B$ un anillo local noetheriano regular y $\mathfrak J$ un ideal de $B$. Para que $A=B/\mathfrak J$ sea un anillo de intersección completa, es necesario y suficiente que el ideal $\mathfrak J$ sea regular (es decir, (16.9.7), que esté engendrado por una sucesión regular de elementos de $B$).
\end{proposition}

Se puede escribir $\widehat A=\widehat B/\mathfrak J\widehat B$ y, puesto que $\widehat B$ es un $B$-módulo fielmente plano y $B$ es noetheriano, $\mathfrak J\widehat B$ es regular si y solo si $\mathfrak J$ lo es (19.1.5, (ii)). Esto muestra (en virtud de ($0$, 17.1.5)), por una parte, que la condición del enunciado es suficiente y, por otra, que para demostrar que es necesaria se puede restringir al caso en que $A$ y $B$ son completos. Por hipótesis, existe entonces un anillo local noetheriano completo regular $B'$ tal que $A$ es isomorfo a un cociente $B'/\mathfrak J'$, donde $\mathfrak J'$ es un ideal regular de $B'$. Consideremos entonces el \emph{producto fibrado} $B\times_A B'=B''$ para los homomorfismos sobreyectivos $B\to A$, $B'\to A$ ($0$, 18.1.2); demostraremos primero el lema siguiente:

\begin{lemma}[19.3.2.1]
\label{IV.19.3.2.1}
Sean $A$, $E$, $F$ tres anillos, $f:E\to A$, $g:F\to A$ dos homomorfismos, y pongamos $G=E\times_A F$ ($0$, 18.1.2).
\begin{enumerate}
  \item[(i)] Si $f$ es sobreyectivo y $E$ y $F$ son anillos noetherianos, lo mismo ocurre con $G$.
  \item[(ii)] Si $A$, $E$ y $F$ son anillos locales y $f$ y $g$ son homomorfismos locales, $G$ es un anillo local y los homomorfismos canónicos $G\to E$, $G\to F$ son locales.
  \item[(iii)] Si se verifican las condiciones de (i) y (ii), $g$ es sobreyectivo y $E$ y $F$ son anillos locales noetherianos completos, lo mismo ocurre con $G$.
\end{enumerate}
\end{lemma}

\begin{enumerate}
  \item[(i)] Teniendo en cuenta ($0$, 18.1.3 y 18.1.5), $G$ es un $A$-anillo aumentado sobre $F$, cuyo ideal de aumento $\mathfrak J'$ se identifica canónicamente con el núcleo $\mathfrak J$ de $f$, y la estructura de $G$-módulo sobre $\mathfrak J'$ se identifica con la de $E$-módulo sobre $\mathfrak J$ por medio del homomorfismo canónico $G\to E$. Si $\mathfrak a$ es un ideal cualquiera de $G$, $\mathfrak a/(\mathfrak a\cap\mathfrak J')$ es isomorfo a $(\mathfrak a+\mathfrak J')/\mathfrak J'$, y por tanto a un ideal de $F$, de modo que es de tipo finito; por otra parte, $\mathfrak a\cap\mathfrak J'$ es isomorfo a un ideal de $E$ contenido en $\mathfrak J$, y por tanto es también de tipo finito, lo que demuestra que $\mathfrak a$ es de tipo finito y, por consiguiente, que $G$ es noetheriano.
  \item[(ii)] Sean $\mathfrak r$, $\mathfrak s$ los ideales maximales de $E$ y $F$, respectivamente; el conjunto $\mathfrak m=\mathfrak r\times_A\mathfrak s$ es evidentemente un ideal de $G$, imagen inversa de $\mathfrak r$ por la proyección $G\to E$ y de $\mathfrak s$ por la proyección $G\to F$. Se deduce inmediatamente que los elementos de $\mathfrak m$ no son invertibles en $G$. Por otra parte, si $(x,y)\notin\mathfrak m$ y, por ejemplo, $x\notin\mathfrak r$, $f(x)$ no pertenece al ideal maximal de $A$; por tanto, puesto que $g(y)=f(x)$ y $g$ es local, se tiene también $y\notin\mathfrak s$; se

\oldpage[IV-4]{195}

concluye que $x$ e $y$ son invertibles, y por tanto también lo es $(x,y)$, lo que termina de demostrar (ii).
  \item[(iii)] La imagen de $\mathfrak m$ por el homomorfismo sobreyectivo $G\to E$ es el ideal maximal $\mathfrak r$ de $E$, por lo que la imagen de $\mathfrak m^n\cap\mathfrak J'$ es $\mathfrak r^n\cap\mathfrak J$; como $\mathfrak J$ es cerrado (y por tanto completo) para la topología inducida por la topología $\mathfrak r$-ádica de $E$ ($0_{\mathrm I}$, 7.3.5), $\mathfrak J'$ es completo para la topología inducida por la topología $\mathfrak m$-preádica de $G$. Por otra parte, $G/\mathfrak J'$, provisto de la topología $\mathfrak m$-preádica, es isomorfo a $F$ provisto de la topología $\mathfrak s$-ádica, y es, por tanto, completo para la topología cociente de la topología $\mathfrak m$-preádica de $G$. Se deduce inmediatamente que $G$ es completo para la topología $\mathfrak m$-preádica.
\end{enumerate}

Establecido este lema, del teorema de Cohen ($0$, 19.8.8) se sigue que $B''$ es isomorfo a un cociente de un anillo local noetheriano completo \emph{regular} $C$. Se deduce entonces de (19.1.2) que la inmersión $\Spec(B')\to\Spec(C)$ es regular; como, por hipótesis, lo mismo ocurre con la inmersión $\Spec(A)\to\Spec(B')$, se concluye de (19.1.5, (iii)) que la inmersión $\Spec(A)\to\Spec(C)$ es regular; pero esta inmersión se escribe también como la compuesta $\Spec(A)\to\Spec(B)\to\Spec(C)$. Ahora bien, como $B$ es regular por hipótesis, la inmersión $\Spec(B)\to\Spec(C)$ es regular (19.1.2); por tanto, también lo es $\Spec(A)\to\Spec(B)$ por (19.1.5, (iv)).

\begin{corollary}[19.3.3]
\label{IV.19.3.3}
Sea $X$ un preesquema localmente noetheriano localmente inmersible en un preesquema regular. Entonces el conjunto de los $x\in X$ en los que $X$ es una intersección completa es abierto en $X$ y contiene al conjunto de los puntos regulares de $X$.
\end{corollary}

En efecto, se puede restringir al caso en que $X=\Spec(B/\mathfrak J)$, donde $B$ es un anillo regular y $\mathfrak J$ un ideal de $B$. El conjunto de los $\mathfrak p\supset\mathfrak J$ de $\Spec(B)$ tales que $\mathfrak J_{\mathfrak p}$ es regular en $B_{\mathfrak p}$ es entonces abierto (16.9.5), y es el conjunto de los puntos de $\Spec(B)$ en los que $X$ es una intersección completa en virtud de (19.3.2).

\begin{corollary}[19.3.4]
\label{IV.19.3.4}
Sean $k$ un cuerpo, $X$ un preesquema localmente de tipo finito sobre $k$, $k'$ una extensión de $k$ y $X'=X\otimes_k k'$. Sean $x'$ un punto de $X'$ y $x$ su proyección en $X$. Para que $X$ sea una intersección completa en $x$, es necesario y suficiente que $X'$ sea una intersección completa en $x'$.
\end{corollary}

Como la cuestión es local sobre $X$, se puede suponer que $X=\Spec(A)$, donde $A$ es una $k$-álgebra de tipo finito y, por tanto, cociente de un álgebra de polinomios $k[T_1,\ldots,T_n]$, de modo que $X$ es un subpreesquema cerrado del esquema regular $Y=\Spec(k[T_1,\ldots,T_n])$ ($0$, 17.3.7). Si se pone $Y'=Y\otimes_k k'=\Spec(k'[T_1,\ldots,T_n])$, $Y'$ es también un esquema regular. Ahora bien, puesto que $\Spec(k')$ es fielmente plano sobre $\Spec(k)$, se deduce de (19.1.5, (ii)) y del hecho de que se trata de preesquemas noetherianos que la inmersión $X\to Y$ es regular en el punto $x$ si y solo si la inmersión $X'\to Y'$ lo es en el punto $x'$. La conclusión se sigue, pues, del criterio (19.3.2).

\begin{corollary}[19.3.5]
\label{IV.19.3.5}
Sean $A$, $C$ dos anillos locales noetherianos, $\rho:A\to C$ un homomorfismo local, $B=C/\mathfrak J$ un anillo cociente de $C$ y $k$ el cuerpo residual de $A$. Se supone que $\rho$ hace de $C$ un $A$-módulo plano y que el anillo $C\otimes_A k$ es regular. Entonces las condiciones siguientes son equivalentes:
\begin{enumerate}
  \item[a)] El ideal $\mathfrak J$ es transversalmente regular relativamente a $A$.
  \item[b)] $B$ es un $A$-módulo plano y $B\otimes_A k$ es un anillo de intersección completa.
\end{enumerate}
\end{corollary}

\oldpage[IV-4]{196}

En virtud de (19.2.4), la condición $a)$ equivale a decir que $B$ es un $A$-módulo plano y que $\mathfrak J\otimes_A k$ (que se identifica con un ideal de $C\otimes_A k$ en virtud de la planitud de $B$ sobre $A$ ($0_{\mathrm I}$, 6.1.2)) es regular en este anillo. Como el anillo $C\otimes_A k$ es regular, en virtud de (19.3.2) equivale a decir (cuando $B$ es un $A$-módulo plano) que $\mathfrak J\otimes_A k$ es un ideal regular de $C\otimes_A k$, o que $B\otimes_A k=(C\otimes_A k)/(\mathfrak J\otimes_A k)$ es un anillo de intersección completa; de donde se sigue el corolario.

\begin{definition}[19.3.6]
\label{IV.19.3.6}
Sea $f:X\to S$ un morfismo plano localmente de presentación finita. Se dice que $X$ es una \emph{intersección completa relativamente a $S$ en un punto} $x$ si la fibra $f^{-1}(f(x))$ es una intersección completa (absoluta) en el punto $x$. Se dice que $X$ es una \emph{intersección completa relativamente a $S$}, o que el morfismo $f$ es un \emph{morfismo de intersección completa}, si $X$ es una intersección completa relativamente a $S$ en cada uno de sus puntos.
\end{definition}

\begin{proposition}[19.3.7]
\label{IV.19.3.7}
Sean $g:X\to S$, $h:Y\to S$ dos morfismos planos localmente de presentación finita, $f:Y\to X$ una $S$-inmersión, $y$ un punto de $Y$, $x=f(y)$ y $s=g(x)=h(y)$. Se supone que la fibra $X_s=g^{-1}(s)$ es regular en el punto $x$. Entonces, para que $f$ sea una inmersión transversalmente regular en el punto $y$, es necesario y suficiente que $Y$ sea una intersección completa relativamente a $S$ en el punto $y$.
\end{proposition}

Decir que $f$ es transversalmente regular en el punto $y$ equivale, en virtud de (19.2.4), a decir que $f_s:Y_s\to X_s$ es una inmersión regular en el punto $y$. Pero, como $X_s$ es regular en el punto $x=f(y)$, esto equivale, en virtud de (19.3.2), a decir que $Y_s$ es una intersección completa (absoluta) en el punto $y$, de donde se sigue la proposición.

\begin{corollary}[19.3.8]
\label{IV.19.3.8}
Sea $f:X\to S$ un morfismo plano localmente de presentación finita. El conjunto $U$ de los $x\in X$ en los que $X$ es una intersección completa relativamente a $S$ es abierto en $X$. Si además $f$ es propio, el conjunto $E$ de los $s\in S$ tales que $X_s=f^{-1}(s)$ es una intersección completa en cada uno de sus puntos es una parte abierta de $S$.
\end{corollary}

Como $E=S-f(X-U)$,\footnote{El texto impreso dice $Y-f(X-U)$; el conjunto ambiente en el que acaba de definirse $E$ es $S$.} la segunda afirmación se sigue de la primera y del hecho de que $f$ es una aplicación cerrada. La primera afirmación es local sobre $X$, por lo que se puede suponer que $X$ es un subpreesquema cerrado de un preesquema de la forma $Z=S[T_1,\ldots,T_n]$ ($T_i$ indeterminadas). Como este último es liso sobre $S$ (17.3.8), sus fibras $Z_s$ son regulares y, por tanto, en virtud de (19.3.7), equivale a decir que $x\in U$ o que la inmersión $X\to Z$ es transversalmente regular en el punto $x$; se concluye, pues, el corolario por (19.2.2).

\begin{proposition}[19.3.9]
\label{IV.19.3.9}
\begin{enumerate}
  \item[(i)] Toda inmersión abierta es un morfismo de intersección completa.
  \item[(ii)] Sea $f:X\to S$ un morfismo plano y localmente de presentación finita que es un morfismo de intersección completa. Para todo cambio de base $g:S'\to S$, $f'=f_{(S')}:X_{(S')}\to S'$ es un morfismo de intersección completa. El recíproco es cierto si $g$ es fielmente plano y cuasicompacto.
  \item[(iii)] Si $f:X\to Y$ y $g:Y\to Z$ son dos morfismos planos localmente de presentación finita que son morfismos de intersección completa, lo mismo ocurre con $g\circ f:X\to Z$.
\end{enumerate}
\end{proposition}

La afirmación (i) se deduce inmediatamente de la definición (19.3.1). Para demostrar (ii), observemos primero que, si $f$ es plano y localmente de presentación finita, lo mismo ocurre

\oldpage[IV-4]{197}

con $f'$ (2.1.4), y el recíproco es cierto si $g$ es fielmente plano ((2.5.1) y (2.7.1)); además, para todo $s'\in S'$, si $s=g(s')$, equivale a decir que $f^{-1}(s)$ es una intersección completa o que $f'^{-1}(s')$ es una intersección completa (19.3.4); de aquí se sigue (ii), teniendo en cuenta que $g$ es sobreyectivo cuando es fielmente plano.

Por último, bajo las hipótesis de (iii), $g\circ f$ es plano y localmente de presentación finita. Por definición (19.3.6), nos reducimos entonces al caso en que $Z$ es el espectro de un cuerpo, y a demostrar que los anillos locales $\sh{O}_{X,x}$ son anillos de intersección completa, lo cual está contenido en el resultado más general siguiente:

\begin{corollary}[19.3.10]
\label{IV.19.3.10}
Sean $f:X\to S$ un morfismo plano localmente de presentación finita que es un morfismo de intersección completa, $x$ un punto de $X$ y $s=f(x)$. Si $\sh{O}_{S,s}$ es un anillo noetheriano de intersección completa, lo mismo ocurre con $\sh{O}_{X,x}$.
\end{corollary}

Evidentemente, se puede restringir al caso en que $S=\Spec(A)$ con $A=\sh{O}_{S,s}$. Demostremos, además, que se puede restringir al caso en que el anillo local $A$ es \emph{completo}. En efecto, pongamos $A'=\widehat A$, $X'=X\otimes_A A'$, y sea $s'$ el único punto cerrado de $S'=\Spec(A')$, que es el único punto de $S'$ sobre $s$; si $f'=f_{(S')}:X'\to S'$, la fibra $f'^{-1}(s')$ es canónicamente isomorfa a $f^{-1}(s)$, puesto que $A$ y $A'$ tienen el mismo cuerpo residual (I, 3.6.4); hay, pues, un único punto $x'\in X'$ sobre $x$ y $s'$, y por consiguiente se tiene $\sh{O}_{X',x'}=\sh{O}_{X,x}\otimes_A A'$. Se deduce que los anillos locales $\sh{O}_{X,x}$ y $\sh{O}_{X',x'}$ tienen \emph{completaciones isomorfas}, pues, en general, si $E$ es un anillo local que es una $A$-álgebra (siendo local el homomorfismo $A\to E$), la completación separada $(E\otimes_A\widehat A)^{\wedge}$ del anillo $E\otimes_A\widehat A$ provisto de la topología producto tensorial es isomorfa a la completación separada de $\widehat E\widehat\otimes_A\widehat A\cong\widehat E$, y por tanto a la completación separada $\widehat E$ de $E$ ($0_{\mathrm I}$, 7.7.1). En virtud de la definición (19.3.1), equivale, pues, a decir que $\sh{O}_{X,x}$ es un anillo de intersección completa o que $\sh{O}_{X',x'}$ es un anillo de intersección completa.

Supongamos, pues, que $A$ es \emph{completo}; se puede suponer además que $X=\Spec(B)$, siendo $B$ cociente de un álgebra de polinomios $C=A[T_1,\ldots,T_r]$. Como los anillos de polinomios sobre un cuerpo son anillos regulares ($0$, 17.3.7), se deduce de (19.3.7) que se puede suponer que la inmersión $X\to Z=\Spec(C)$ es transversalmente regular relativamente a $S$; por tanto, se puede suponer que $B=C/\mathfrak J$, donde $\mathfrak J$ está engendrado por una sucesión regular $(g_i)_{1\leq i\leq n}$ de elementos de $C$. Por otra parte, puesto que $A$ es completo, se puede, por hipótesis, escribirlo como $A'/\mathfrak R$, donde $A'$ es un anillo local regular y $\mathfrak R$ un ideal de $A'$ ($0$, 19.8.8), y la hipótesis de que $A$ es un anillo de intersección completa implica que $\mathfrak R$ está engendrado por una sucesión regular $(f_j)_{1\leq j\leq m}$ (19.3.2). En el anillo de polinomios $C'=A'[T_1,\ldots,T_r]$, los elementos $f_j$ ($1\leq j\leq m$) siguen formando una sucesión regular que engendra el ideal $\mathfrak R'=\mathfrak R[T_1,\ldots,T_r]$ ($0$, 15.1.4); como $C'/\mathfrak R'=C$, se ve que, si para cada $i$, $g'_i$ es un elemento de $C'$ de imagen $g_i$ en $C$, la sucesión formada por los $f_j$ ($1\leq j\leq m$) y los $g'_i$ ($1\leq i\leq n$) es regular en $C'$; si $\mathfrak J'$ es el ideal de $C'$ que engendra, se tiene $B=C'/\mathfrak J'$, de donde se sigue la conclusión en virtud de (19.3.2), puesto que $C'$ es un anillo regular ($0$, 17.3.7).
